% V20 recollation: French physical 66--69 (67/68 duplicate); English 44--45.
% Scope ends before Proposition 6.1.  Supplementary arguments are in the audit.
\source{法文原稿物理第66页“命题2.1”起，续第67、68页的同一稿页，至第69页上部5.7止；英文编辑稿物理第44--45页}

命题2.1可以不作改动地移用到这里。在命题2.2中，应当读作：充要条件是，对每个$x\in Z$都有\cmt{V20-EXT001}{法第66页在$x$上确有闭包横线；英文第44页漏作$\dim x$。这里是点在$P$中闭包的维数，不是局部环的维数；对任意集合$Z$，也不能换成$\dim\overline Z\leq m-1$。英文把\emph{se transpose}译成“translates (?)”，法文这里只说原命题可移用。点闭包判据及可构造性限制的译校论证另见本版核查报告。}
\[
 \dim\overline{\{x\}}\leq m-1.
\]

为作证明，例如可以把余维$m$的一般线性子簇看成$m$个独立的一般超平面的交，从而归约到2.6。让迪厄多内自己想办法吧。\cmt{V20-EXT002}{原稿接写\emph{Dieudonn\'e demerdetur.}，这是作者对编写者的插语。英文先标作拉丁语，又附译者疑问：“或者这是稍显粗俗的法语？”这里译出其意，保留原句及这项英译附注，不把它当作证明步骤。“独立的一般超平面”涉及联合参数的一般点，不能解释为任取若干分别满足某个好性质的超平面。}

从编写的角度说，如果从一开始就取一般的$m$（我觉得这样更好），那么似乎最好同时证明2.6；其中当然把
\[
 \dim X-1\quad\text{替换为}\quad\dim X-m,
\]
并约定公式中维数小于零表示所讨论的集合为空。\cmt{V20-EXT003}{原文是对2.6的修改指示，仍须保留被引命题的各分支像维数条件。像维数小于$m$的分支可能被一般截面完全避开，不能把这种分支的维数照搬入非空截面的维数公式。本段没有给出任意线性截面均有预期维数的断言。}

推论2.3中，把“有限”替换为“维数不超过$m-1$”，即可得到相应表述。推论2.4也一样。2.5同样处理：把
\[
 \dim f(X_i)>0\quad\text{替换为}\quad\dim f(X_i)\geq m;
\]
2.7中也作同样的替换。\cmt{V20-EXT004}{本段及下文沿用原稿的像维数简记$\dim f(W)$，按前文约定指$\dim\overline{f(W)}$。2.7的原对象是$F$的关联素循环$Z_i$，故该处修改的是$\dim f(Z_i)$的下界，不是把$Z_i$改名为$X_i$。各项“有限”的替换须留在各自原条目的像集、坏点集和有限性假设中理解。}

命题2.8照原样仍然成立；在2.9中，把“有限”替换为“维数不超过$m-1$”。2.10、2.11也一样。

陈述2.12照原样仍然成立，证明基本不变（也参见下文关于第8节的说明）。在2.13中，把“有限”替换为“维数不超过$m-1$”。定理2.14照原样仍然成立；2.15则作上述“有限”的替换。2.16照原样有效；2.17也把“有限”替换为“维数不超过$m-1$”。2.18照旧。\cmt{V20-EXT005}{这句跨越法第66页与第67/68页；第67、68页从同一续句开始，是同一稿页的两次扫描，不是两份应相继翻译的文字。第8节前的手写\emph{plus bas}已译作“下文”。“证明基本不变”是作者的推广判断；例如2.12原有的无分歧假设并未在这里取消。本次不据这些逐项指示声称各条推广均已完成独立证明。}

对于3.1，可以在假设
\[
 \dim f(X)\geq m+1
\]
的条件下，对任意$m$作出相应陈述。不过，我建议主要陈述仍保留在超曲面情形，而把一般情形作为推论或附记给出（用通常的办法，取独立的一般超平面，就能立即导出它）。除非觉得把引理3.1.1的一个推广版本明确写出来也颇有意思……\cmt{V20-EXT006}{法文为\emph{\`A moins qu'il soit amusant d'expliciter\ldots}，是对前一编写建议保留另一种可能；英文译成“At least”，旧译据之作“至少”，改变了句间关系，现据法稿改正。省略号及尚待写出的引理均保留，不补造引理的推广陈述。}

对于3.2，条件应读作$\dim f(X)\geq m+1$。在3.3中，把$\dim f(X_i)\geq2$替换为$\dim f(X_i)\geq m+1$；在$\sigma$的定义中，把$\dim f(Z)=0$替换为$\dim f(Z)\leq m-1$。\cmt{V20-EXT007}{法第16页3.3与第67/68页在此使用相同的手写闭集族记号；英文第13页转写为$\sigma$，第45页却转写为$G$。正文沿用前文$\sigma$，指3.3中的闭集族，与本节的泛秩$m$商模$G$区别。这里的两项改动分别作用于不可约分支的像与闭集族的定义。}

第4节的一般讨论可照原样用于任意$m$。4.2和4.3也一样；其中，4.3的b）在(v)、(vi)两种情形中的维数条件改为
\[
 \dim f(X)\geq m+1.
\]
4.4的b）作类似修改。\cmt{V20-EXT008}{法第67、68页合写“4.2和4.3”，随后指向b）的(v)、(vi)；核法第18、19页，4.2的旧(vi)已划销，现行的对应维数条件在4.3(b)。阈值的手改笔画在两次扫描中均有叠压，英译保留“$m-1$ or $m+1$ (illegible, ask AG)”；本版仍将字形层面的歧义标明，不声称仅凭图像已彻底消歧。正文采用可由被引条目确定的$m+1$：$m=1$须回到4.3(b)的下界2，同页3.1、3.2的推广也给此下界；$m-1$的读法被光滑平面二次曲线的一般直线截面反驳，详见核查报告。4.4中的改动针对各分支的像维数；其$(C_j)$中的连通性指标与本节余维$m$须分开，作者在原4.4后所加的拟有限限制或像维数版本也仍须保留。}

5.1中的那段议论照旧。\cmt{V20-EXT009}{法第67/68页清楚写\emph{Le la\"ius 5.1. passe tel quel.}；英文保留“[Le laius 5.1 Fr or Latin 67]”。这是对语言和稿页的编辑标记，67不是数学条目号。}
在5.2中，须记住：$\phi$现在成为$G_T(1)$的一个截面$\phi^{(m)}$，其中
\[
 T=X\times_S\operatorname{Gr}^{m}.
\]
对于$\xi\in\operatorname{Gr}^{m}$，它诱导截面\cmt{V20-EXT010}{法第67/68页此处有数处划销的措辞和旧对象，改后保留$T=X\times_S\operatorname{Gr}^m$及所示纤维模。$G(\xi)$是本节泛商模在$\xi$的$m$维纤维，张量式将它拉回到$X_{k(\xi)}$；它不是第5节曾用$G$表示的截面限制模。英文所排重复的$k(\xi)$下标在此按基变换展开。无法从划销层可靠复原的旧文字不补入定稿，具体位置存于核查报告。}
\[
 \phi^{(m)}_\xi\in
 \Gamma\bigl(X_{k(\xi)},
 \mathcal O_{X_{k(\xi)}}(1)\otimes_{k(\xi)}G(\xi)\bigr).
\]


一般而言，我们将在第19段说明：给定一个预概形上秩$m$的局部自由模的截面$\phi$，以及一个模$F$，怎样称这个截面为$F$-正则。用一组局部基来表达，就是$\mathcal O_X$的$m$个截面组成一个$F$-正则序列（还须验证这不依赖于所选的基）。\cmt{V20-EXT011}{法第67、68页两处均写$m$；英文将前一处“秩$m$”误排为“秩$n$”，这是英译转写错，不是法稿秩数不一致。括号跨到法第69页首行结束。原稿仍要求验证基无关性；本版仅在核查报告中另证诺特局部环、有限非零模及截面零点处的相应结论，并说明Koszul正则性与逐项非零因子条件在一般环境、特别是零点集外的区别。}
当$m=1$时，就有5.2中提到的那个显然的内在解释。采用这一术语约定，5.3照原样仍然成立，5.4也是如此。

附注5.5的第一部分可以推广到任意$m$的情形：若$F_s$满足$(S_m)$，则关于$\phi_\xi$的上述正则性条件，可以用纯粹的维数条件来表达。\cmt{V20-EXT012}{法第69页为小写下标$s$，条件$(S_m)$为手写修订；英文排作$F_S$。这里承接5.5，$s\in S$是$\xi$的像，正则性须在$X_{k(\xi)}$上对$F_{k(\xi)}$检验，不能改成总空间上的$F$-正则性。原文没有写出维数判据的公式；报告中给出的局部支集余维判据是注明假设的译校补证，未将$(S_m)$误读为任意截面均正则。}

附注5.5的第二部分对于任意$m$都照原样有效。定理5.6也照原样推广，5.7也是如此。\cmt{V20-EXT013}{这些仍是原作者的推广指示。5.5第二部分关于Cohen--Macaulay商的结论须保留截面正则性；5.6的有限表示与相对$S$平坦、5.7的$\Supp F$在$S$上固有等原假设也未被删除。本次没有凭“照原样”补写并宣称验证全部推广证明，亦没有将第5节的性质清单及其引用问题视为在此自动消失。}
% End of owned recollation range; next source sentence is Proposition 6.1.
